\section{The Optimization Pipeline Under Study}
\label{sec:subject}
% NEW (written for the TOMPECS version; not in the original paper)
This section describes the system whose behaviour the rest of the paper evaluates: the loop
itself, the profiling evidence it can place in the prompt, and the structure of that prompt.

\subsection{Pipeline}
The framework takes a source file, a build command and a run command, and closes the loop
\emph{profile $\rightarrow$ prompt $\rightarrow$ generate $\rightarrow$ build $\rightarrow$ verify
$\rightarrow$ measure}, with bounded repair. The optimized file returned by the model replaces the
original in a private copy of the application tree, so the build, the callers and the input data are
those of the real program. If the build or the correctness check fails, the error summary is appended
to the prompt and the model is asked to repair; we allow three iterations before recording the run as
a failure.

\subsection{Profiling data}
We profile with Nsight Compute in two passes: an uncapped pass over every launch that yields true
launch counts and per-launch durations, and a capped \texttt{-{}-set full} pass that yields the detail
pages. From the resulting reports we derive three kinds of prompt context:
\begin{itemize}\itemsep2pt
  \item \textbf{PC sampling} (\texttt{-{}-page source}): per-source-line warp stall reasons, which we
    aggregate to the lines with the highest stall ratios.
  \item \textbf{Roofline} (\texttt{-{}-page details}): per-kernel throughput relative to DRAM and compute
    peaks, plus the analysis rules Nsight Compute reports.
  \item \textbf{Importance analysis}: per-kernel aggregate cost (launch count $\times$ per-launch time)
    used to rank which kernel is worth editing at all.
\end{itemize}

\subsection{Prompt structure}
The prompt has five regions: (i) the source file, (ii) the profile summaries above, (iii)
hardware-specific guardrails (target architecture, shared-memory limits, unavailable intrinsics), (iv)
general pitfall-avoidance guardrails (entry-point preservation, no signature changes, no new top-level
symbols, semantic preservation), and (v) an output-format specification requiring a single fenced code
block plus structured lists of applied and rejected optimizations. Section~\ref{sec:ablation} ablates
these regions; Section~\ref{sec:profabl} removes region (ii) entirely.
